\textbf{Detector láser para riesgo de hidrógeno: selección y restricciones INC-24.}\label{sec:sd4-inc24}
El producto VLC-EX-US documentado por Xtralis emplea tecnología láser y declara aplicación Class I Division 2, grupos A/B/C/D. Constituye una referencia técnica más específica que VLF-500-00-UL para una clasificación que incluya hidrógeno, pero no se adjudica automáticamente a F04/F05: la página actual también publica fin de vida de VLP/VLC, y la familia no sustituye certificado vigente de la variante, condiciones especiales, muestreo, escape, accesorios ni confirmación independiente. Synaptix conserva cuatro VLF-500-00-UL seleccionados en zonas ordinarias y dos VEP-A00-1P-UL seleccionados en INC-30, con clasificación, condiciones y confirmación pendientes antes de compra/energización. No se utiliza VLX-100 por analogía entre nombres Ex ni se incorpora VLI como equivalente láser. El expediente del banco debe identificar emisión de las tres condiciones normal, falla de extracción y posaislamiento, alcance del área peligrosa y aptitud del conjunto; su confirmador independiente debe tener certificación y conexión compatibles con el panel Cheetah. El cumplimiento integral permanece pendiente de el expediente íntegro de la selección actual y esos documentos \parencite{l24-xtralis-vlc}.

\textbf{Recipientes pequeños de cohorte EN EXT-30.}\label{sec:sd4-ext24}
Se sustituyen los tres 70-264 DOT y sus reemplazos por 70-350 EN de 16 L del manual 06-542, llenados 11/11/15 kg por cohorte, seis pequeños en total. Diámetro 229 mm, altura 635 mm, tara 16 kg y masa conjunta con agente 170 kg. Los llenados tienen m/V de dos tercios; las concentraciones aproximadas y condiciones por riesgo se desarrollan en EXT-30. No se trasladan las cotas DOT de 178 mm/857,3 mm/14,7 kg al nuevo recipiente. El gabinete conserva su reserva y exige conexiones, control térmico y anclaje verificados; hidráulica/listado/retención y seguridad gas siguen pendientes.

\textbf{Plano y cálculo completo de portátiles PORT-24.}\label{sec:sd4-port24}
Se seleccionan veintinueve equipos activos: veintiséis Amerex B441 ABC 4A:80B:C, un Amerex 398 Halotron I 2A:10B:C y dos Amerex 322 CO$_2$ 5B:C. Cinco equipos de reserva local, tres B441 más uno de cada otro modelo, conservan al menos el 10\,\% por tipo redondeado hacia arriba: treinta y cuatro equipos totales. P01--P09 mantienen modelos y ubicación interior de PORT-21; P10--P29 son B441 y se fijan en el plano exterior actualizado. El potencial chileno acreditado se recibe al habilitar y no se deduce de UL. Las obligaciones de certificación, servicio técnico certificado, inspección mensual, mantenimiento anual, reemplazo antes de retirada, capacitación y bitácora de PORT-21 siguen vigentes \parencites[pp.~9, 11--12]{l21-amerex}[arts.~45--51]{l21-ds594}[arts.~23--33]{l21-ds44}.

La reserva nueva ocupa $[-2;32]\times[-2;37]$ m, 34$\times$39 m, sin atribuir esas dimensiones al edificio existente. El módulo conserva origen interior $(10;8)$ y 10$\times$7 m, con reserva exterior propia de muros de 0,50 m por lado, coordinada con blindaje, bastidor y aislamiento: obstáculo exterior de 11 por 8 m. Las UTA inferiores pasan a $[7,5;10,5]\times[20;22]$ y $[19,5;22,5]\times[20;22]$; las superiores, a $[6;9]\times[31;33]$ y $[21;24]\times[31;33]$. Los gabinetes de recipientes ocupan $[11;14]\times[31;33]$ y $[16;19]\times[31;33]$. Agua A/B pasa a $[7,5;9,5]\times[4;6]$ y $[20,5;22,5]\times[4;6]$. Grupos, cubetos, bancadas, frío y trabajo/stock conservan las demás cotas del plano. Esas relocalizaciones eliminan pasos originales de 0,5--1,0 m entre reservas. Cada eje de ruta se aleja al menos 0,75 m de las reservas de máquinas y cubetos: el corredor continuo asignado es 1,50 m, con barreras, puertas y nichos fuera de él. La cota propia del corredor no acredita holguras normativas de máquinas todavía sin variante; si la variante excede su reserva, se redimensiona el conjunto antes de compra.

\begin{longtable}{|p{0.12\linewidth}|p{0.67\linewidth}|}
\hline \textbf{Posición} & \textbf{Coordenadas de toma en metros, origen del plano exterior} \\\hline
\endfirsthead
\hline \textbf{Posición} & \textbf{Coordenadas de toma en metros, origen del plano exterior} \\\hline
\endhead
P10 & $(3;17,25)$ \\\hline
P11 & $(27;17,25)$ \\\hline
P12 & $(4,75;11,5)$ \\\hline
P13 & $(25,25;11,5)$ \\\hline
P14 & $(4,75;3,5)$ \\\hline
P15 & $(25,25;3,5)$ \\\hline
P16 & $(12,25;26)$ \\\hline
P17 & $(17,75;26)$ \\\hline
P18 & $(16;3,75)$ \\\hline
P19 & $(3;22,75)$ \\\hline
P20 & $(27;22,75)$ \\\hline
P21 & $(-0,75;3,5)$ \\\hline
P22 & $(30,75;3,5)$ \\\hline
P23 & $(7,5;30,25)$ \\\hline
P24 & $(22,5;30,25)$ \\\hline
P25 & $(15;32)$ \\\hline
P26 & $(11,25;20)$ \\\hline
P27 & $(18,75;20)$ \\\hline
P28 & $(7,5;36,25)$ \\\hline
P29 & $(22,5;36,25)$ \\\hline
\end{longtable}
Las coordenadas representan el punto accesible de retiro; el poste/nicho se sitúa hacia el lado opuesto de la franja libre de paso, con planta propia de 0,35$\times$0,25 m, sin invadir huella de máquina ni eje. P19/P20 se añaden al norte de grupos; P21/P22, al lado exterior de bancadas; P23/P24 protegen UTA superiores; P25 atiende recipientes; P26/P27 protegen climatización TI y UTA inferiores. P28/P29 protegen las estaciones de impulsión superiores desde el norte, justificando dos B441 adicionales. P16/P17 se desplazan para respetar 0,75 m de separación a plataformas de frío. Ninguna posición exige ingresar a un banco para retirar un equipo ni atravesar un cubeto.

Se calcula la distancia por rutas ortogonales en retícula de 0,25 m, excluyendo las reservas de máquinas/cubetos ampliadas en 0,75 m. La ampliación incluye la envolvente de equipo y sus zonas asignadas, no sólo el símbolo de catálogo. El grafo sólo admite segmentos de corredor libre; puertas, muro y máquina son obstáculos. Se resuelven distancias a la estación más cercana y se comprueba cada tramo del perímetro, con muestras separadas 0,25 m; una posición intermedia añade como máximo 0,125 m al valor de la muestra cercana porque puede recorrer ese mismo borde. Esa cota continua explica por qué se comprueba más que sus vértices. Trabajo y stock se consideran superficies de ocupación, no obstáculos: se verifica toda su retícula interior y se suma 0,25 m para cualquier posición entre cuatro muestras. No se ofrece protección del pavimento vacío como si fuera combustible; se cubren todos los lugares de trabajo y equipos identificados.

\begin{longtable}{|p{0.30\linewidth}|p{0.17\linewidth}|p{0.24\linewidth}|}
\hline \textbf{Zona protegida} & \textbf{Área de reserva} & \textbf{Cota continua del recorrido} \\\hline
\endfirsthead
\hline \textbf{Zona protegida} & \textbf{Área de reserva} & \textbf{Cota continua del recorrido} \\\hline
\endhead
grupoA & 24 m$^2$ & 6,625 m \\\hline
grupoB & 24 m$^2$ & 6,625 m \\\hline
tanqueA & 12 m$^2$ & 8,875 m \\\hline
tanqueB & 12 m$^2$ & 8,875 m \\\hline
bancadaA & 28 m$^2$ & 7,125 m \\\hline
bancadaB & 28 m$^2$ & 7,125 m \\\hline
frioTI\_A & 4 m$^2$ & 7,125 m \\\hline
frioTI\_B & 4 m$^2$ & 7,125 m \\\hline
frioA & 20,0 m$^2$ & 6,125 m \\\hline
frioB & 20,0 m$^2$ & 6,125 m \\\hline
UTA\_A1 & 6,0 m$^2$ & 6,375 m \\\hline
UTA\_B1 & 6,0 m$^2$ & 6,375 m \\\hline
UTA\_A2 & 6 m$^2$ & 6,625 m \\\hline
UTA\_B2 & 6 m$^2$ & 6,625 m \\\hline
aguaA & 4,0 m$^2$ & 8,875 m \\\hline
aguaB & 4,0 m$^2$ & 7,875 m \\\hline
cilTI & 6 m$^2$ & 5,875 m \\\hline
cilPeq & 6 m$^2$ & 5,875 m \\\hline
imp\_A1 & 2,5 m$^2$ & 6,625 m \\\hline
imp\_B1 & 2,5 m$^2$ & 6,625 m \\\hline
imp\_A2 & 2,5 m$^2$ & 6,625 m \\\hline
imp\_B2 & 2,5 m$^2$ & 6,625 m \\\hline
trabajo & 6 m$^2$ & 7,0 m \\\hline
stock & 3,0 m$^2$ & 6,0 m \\\hline
\end{longtable}
Cada reserva de la tabla tiene menos de 150 m$^2$, atendida por B441 de potencial 4A de catálogo, con cota de traslado máxima 8,875 m, inferior a 9 m. Los cubetos reciben además potencial B de catálogo 80B, superior al mínimo chileno 40B; su certificación local debe acreditar ese mínimo. El cálculo conserva márgenes explícitos y rutas alrededor de los equipos, sin sumar simplemente sus lados ni cruzar sus huellas. La Figura~\ref{fig:recinto-exterior-coordinado} dibuja las rutas críticas y todas las estaciones; las rutas pueden compartirse como corredor de circulación, sin invadir equipos. El expediente de aceptación incluirá plano, inventario y certificados vigentes; la recepción futura comprobará dimensiones reales, recorridos y acceso libre. Esa recepción no se presenta como realizada.

La elección de ruta se aplica para recuperación y evacuación desde un sector seguro. Una persona no rodeará la fuente activa de fuego ni atravesará humo, escape caliente o atmósfera H$_2$ para alcanzar un extintor: si el corredor hacia su estación está afectado, se retira por salida segura y utiliza otra estación sólo desde ese lado, sin intervención dentro del área peligrosa. El portátil es protección secundaria; la evacuación y aislamiento competentes prevalecen sobre combatir un incendio desarrollado.

\textbf{Rutas interiores PORT-24.} El cálculo interior mantiene las seis paredes de sector, permite cruzar sólo las puertas indicadas y excluye cuatro gabinetes de batería, dos racks y dos climatizadores. Se fijan las huellas UPS A $(1,43;0,85)$--$(2,25;1,35)$ y UPS B $(7,75;0,85)$--$(8,57;1,35)$, frente hacia salida oeste/este y servicio por ese lado: no se presupone acceso detrás de la UPS contra pared. En TI, las huellas AR3104 y s-MEXT permanecen según Figura~\emph{Plano de distribución proyectada del recinto técnico y zonas asociadas} (SD4); sus pasillos anteriores/posteriores y el pasillo de cada banco son los espacios ocupables. Se calcula una retícula de 0,05 m con eje a al menos 0,60 m del equipo/muro, corredor útil propio de 1,20 m; la descarga/retirada del extintor mural añade el desplazamiento al punto de toma. La cota para cualquier posición entre muestras es 0,05 m. F01: máximo 6,55 m hasta P01, con rodeo por el pasillo de racks; F02/F03: 2,90 m hasta P02/P03, cruzando sólo salida lateral de cada UPS; F04/F05: 4,70 m hasta P04/P05, por su pasillo hacia puerta norte; F06: 3,50 m hasta P06, sin atravesar el tabique TI. Incluyen 0,50 m de aproximación a P01; 1,95 m desde eje interior UPS por puerta y lado exterior hasta toma; 0,85 m desde eje norte de bancos y 0,40 m hasta P06. Los pasos de servicio no admiten acopio ni nueva huella. El proyecto de ocupación distribuye personas exclusivamente en esos corredores y zonas de trabajo; no se acredita circulación detrás de gabinetes ni a través de equipos. Cada sector ocupa menos de 150 m$^2$ y dispone de B441 propio 4A:80B:C, sin depender de extintor de otro compartimento. P07 y P08/P09 son complementarios: el Halotron permanece sólo en los 82,5 m$^3$ de TI y el CO$_2$ no reemplaza cobertura A.
